% !TeX spellcheck = en_US
\documentclass[11pt]{article}
\usepackage{exerciseHead}

%%% ENCODING, FONT & LOCALIZATION
\usepackage[utf8]{inputenc} % use unicode input-encoding (allows unicode characters in the .tex files, default: Asci)
\usepackage{lmodern} % vectorized version of the "computer modern" font (default) (https://tex.stackexchange.com/q/1390/148164, https://tex.stackexchange.com/q/147194/148164)
\usepackage[ngerman, english]{babel} %% last language is main
%\shorthandoff{"} %babel breaks tikz-cd otherwise (activate if used)
\usepackage{csquotes} %% biblatex wants csquotes (localized quotes) if babel is present

%%% GENERAL IMPROVEMENTS
\usepackage[unicode, colorlinks]{hyperref} %% clickable links
\usepackage{enumitem} %% better (customizable) enumerations (e.g. \begin{enumerate}[label=(\roman*)])


%%% STYLING
% \setlength{\parindent}{0em} % No indentation for new paragraphs
\usepackage{emptypage} % removes page numbers on empty pages
% \makeatletter \@twosidefalse \makeatother
\usepackage{xcolor}
\usepackage[many]{tcolorbox}
\usepackage{fancyhdr}
\usepackage{mdframed}
\usepackage{booktabs}
\usepackage{multirow}

\usepackage{marvosym} % special letters


%%% GRAPHICS & PLOTS

% \usepackage{graphicx} % enable pictures
\usepackage{wrapfig} % text wrapping around figures
\usepackage{wrapstuff}
\usepackage{tikz, tikz-cd, pgfplots} % plots and sketches
\usetikzlibrary{bending,positioning,calc}
\usetikzlibrary{decorations.pathreplacing}
\usetikzlibrary{arrows.meta}
\pgfplotsset{compat=1.18}


\usepackage{subcaption}

%%% CODE
\usepackage[ruled,linesnumbered]{algorithm2e}

% fix algomathdisplay (see: https://tex.stackexchange.com/a/611426/148164)
\makeatletter
\renewenvironment{algomathdisplay}
 {\[}
 {\@endalgocfline\vspace{-\baselineskip}\]{\DontPrintSemicolon\;}}
\makeatother
%%% MATH
\usepackage{mathtools} % fixes some things from amsmath and adds some new features
\usepackage{amsthm, amssymb} % amsthm -> proof env, amssymb -> additional symbols
% \usepackage{thmtools} % restateable theorem environments
\usepackage{aligned-overset} % https://tex.stackexchange.com/questions/257529/overset-and-align-environment-how-to-get-correct-alignment
%% usage: \overset{something}&{=}
\usepackage{xfrac}

%%% Special Symbols
% \usepackage{bbm} % blackboard numbers: usage \mathbbm{1} (amssymb \mathbb{1} looks bad)
% \usepackage{cancel} % strike through equation parts to indicate they cancel out
% %% contradiction symbol
% \usepackage{wasysym} % lightning symbol
% \newcommand*{\contradiction}[1][]{\quad\text{{\LARGE\lightning} #1}\ }

%%% BIBLIOGRAPHY
\usepackage[numbers]{natbib}

% Editing
\usepackage[draft]{fixme}
\fxsetup{theme=color}

\usepackage{xsim}


%% AMS Theorem Environments
\theoremstyle{plain}% default
\newtheorem{proposition}{Proposition}[section]
\newtheorem{lemma}[proposition]{Lemma}
\newtheorem{corollary}[proposition]{Corollary}
\newtheorem{theorem}[proposition]{Theorem}
\newtheorem{fact}[proposition]{Fact}

\theoremstyle{definition}
\newtheorem{definition}[proposition]{Definition}
\newtheorem{example}[proposition]{Example}
\newtheorem{assumption}[proposition]{Assumption}

\theoremstyle{remark}
\newtheorem{remark}[proposition]{Remark}
%%

\newcommand*{\dims}{d}

\DeclarePairedDelimiter{\abs}{\lvert}{\rvert}
\DeclarePairedDelimiter{\Set}{\{}{\}}
\NewDocumentCommand{\Indicator}{s m}{%
  \IfBooleanTF{#1}
    {\mathbf{1}_{\left\{#2\right\}}} % starred: wrap in braces
    {\mathbf{1}_{#2}}                 % unstarred: simple subscript
}

% spaces
\newcommand*{\real}{\mathbb{R}}
\newcommand*{\nat}{\mathbb{N}}
\newcommand*{\integer}{\mathbb{Z}}
\newcommand*{\complex}{\mathbb{C}}
\newcommand*{\rational}{\mathbb{Q}}

% Topology
\newcommand*{\closure}[1]{\overline{#1}}

% Probability
\newcommand{\E}{\mathbb{E}}
\renewcommand{\Pr}{\mathbb{P}}
\DeclareMathOperator{\Var}{Var}
\DeclareMathOperator{\Cov}{Cov}
% \newcommand*{\normal}{\mathcal{N}}
\newcommand{\normal}[1]{{\color{blue}\boldsymbol{\mathcal{N}}}\left(#1\right)}
\newcommand*{\density}{\varphi}
\DeclareMathOperator{\iid}{iid}
\newcommand*{\given}{\mid}

\newcommand*{\Binomial}{\mathrm{Bin}}
\newcommand*{\Poisson}{\mathrm{Poi}}

% random functions
\newcommand*{\rf}{\mathbf{f}}

% Linear Algebra
\DeclareMathOperator{\sgn}{sgn}
\DeclareMathOperator{\Span}{span}
\newcommand{\identity}{\mathbb{I}}
\newcommand*{\transpose}{T}

% Optimization
\DeclareMathOperator*{\argmin}{arg\,min}
\DeclareMathOperator*{\argmax}{arg\,max}

% Machine Learning
\newcommand*{\cost}{J}
\newcommand*{\Cost}{\mathbf{J}}
\newcommand*{\loss}{\ell}
\newcommand*{\noise}{\varsigma}
\newcommand*{\param}{w}
\newcommand*{\model}{\phi}
\newcommand*{\X}{X}
\newcommand*{\Y}{Y}


% Numerical Analysis
\newcommand*{\bigO}{\mathcal{O}}


%% =========== Color (use RGB for digital, CMYK for print) ==========

% ----- Uni Mannheim Corporate Design -----
\definecolor{primary}{RGB}{0,48,86}
% \definecolor{primary}{cmyk}{1,0.6,0.1,0.6}

\colorlet{primary10}{primary!10}
\colorlet{primary25}{primary!25}
\colorlet{primary40}{primary!40}
\colorlet{primary55}{primary!55}
\colorlet{primary70}{primary!70}
\colorlet{primary85}{primary!85}

\definecolor{accent}{RGB}{179,182,185}
% \definecolor{accent}{cmyk}{35,25,25,0}

% ------------------------------------------

% --- convenient color access ---
\newcommand*{\colorPrimary}[2][]{{\color{primary#1} #2}}
\newcommand*{\colorAccent}[1]{{\color{accent} #1}}

 % AFTER the solution true/false macro !!

% \editmodetrue

\ifeditmode
  \NewDocumentEnvironment{solution}{O{Solution} +b}{\begin{proof}[#1] #2\end{proof}}{}
  \newtheorem{exercise}{Exercise}[section]
\else
  \usepackage{xsim}
  \DeclareExerciseEnvironmentTemplate{theorem}{%
    \par\addvspace{10pt}
    \noindent\textbf{%
      \XSIMmixedcase{\GetExerciseName}~\GetExerciseProperty{counter}%
    }{\normalfont%
      \GetExercisePropertyT{subtitle}{ (\PropertyValue)}}\textbf{. }
    \IfInsideSolutionF{%
      \GetExercisePropertyT{points}{%
        \marginpar{%
          \printgoal{\PropertyValue}%
          \GetExercisePropertyT{bonus-points}{+\printgoal{\PropertyValue}}%
          \,\IfExerciseGoalSingularTF{points}
            {\XSIMtranslate{point}}
            {\XSIMtranslate{points}}%
        }%
      }%
    }%
  }{\par\addvspace{\baselineskip}}

  \renewcommand\printsolutions{%
    \def\currentchapter{}%
    \def\currentsection{}%
    \def\lastchapter{}%
    \def\lastsection{}%
    \chapter*{Solutions\addcontentsline{toc}{chapter}{Solutions}}%
    
    \ForEachUsedExerciseByType{%
      \let\lastchapter\currentchapter
      \let\lastsection\currentsection
      \edef\currentchapter{\ExercisePropertyGet{##1}{##2}{chapter-value}}%
      \edef\currentsection{\ExercisePropertyGet{##1}{##2}{section-value}}%
      % \ifx\lastchapter\currentchapter\else
      %   \addcontentsline{toc}{chapter}{Solutions Chapter \ExercisePropertyGet{##1}{##2}{chapter}}
      %   \chapter*{Solutions Chapter \ExercisePropertyGet{##1}{##2}{chapter}}
      % \fi
      \ifx\lastsection\currentsection\else
        \section*{Exercises \ExercisePropertyGet{##1}{##2}{section}}%
        \addcontentsline{toc}{section}{Exercises \ExercisePropertyGet{##1}{##2}{section}}%
      \fi
      \XSIMprint{solution}{##1}{##2}%
    }%
  }
  \xsimsetup{
      exercise/template=theorem,
      solution/template=theorem,
      exercise/the-counter = \thesection.\arabic{exercise},
      exercise/within=section,
      % solution/print=true,
  }
  % \SetExerciseParameters{exercise}{counter=theorem}
\fi
  \xsimsetup{
      exercise/the-counter = \arabic{exercise},
      solution/print=true,
  }


% Meta Information %%%%%%%%%%%%%%%%%%%%%%%%%%%%%%%%%%%

% included in all sheets (-> only one update per semester)
\course{Compl\'ements de probabilit\'es et statistiques}
\semester{Spring 2026}
\author{Felix Benning}

 % (file) for one update per semester
\sheetNumber{1}

%%%% Document %%%%%%%%%%%%%%%%%%%%%%%%%%%%%%%%%%%%%%%%
\begin{document}

\maketitle

\textbf{
The goal is to have something to talk about in the oral exam.  Skip an exercise
if you get stuck. That is why there are many.
% I am more interested in how you
% would solve them than the actual solution, i.e.\ do not prioritize tedious
% calculations.
}


\begin{exercise}[subtitle={Convergence in probability but not \(L^p\)}]
    \label{ex: convergence in probability but not in Lp}
    Let \(U\sim \uniform(0,1)\) be a standard uniform random variable and define
    \[
        X_n \coloneq n \ind{[0, 1/n)}(U)
    \]
    Show that
    \begin{enumerate}[label=(\alph*), noitemsep]
        \item\label{it: (conv in prob but not in Lp) in prob}
        the sequence \((X_n)_{n\in \nat}\) converges to \(0\) in probability,

        \item\label{it: (conv in prob but not in Lp) not in L1}
        the sequence \((X_n)_{n\in \nat}\) does not converge to \(0\) in \(L^1\),

        \item\label{it: (conv in prob but not in Lp) for all q}
        For every \(p \ge 1\) construct a sequence \((Y_n)_{n\in \nat}\) such that
        \(Y_n\xrightarrow{L^q} 0\) for all \(q < p\) but not for \(q \ge p\).
    \end{enumerate}
\end{exercise}
\begin{solution}
\begin{enumerate}[wide=0pt]
    \item[\ref{it: (conv in prob but not in Lp) in prob}]
    We have for any \(\epsilon > 0\)
    \[
        \prob{|X_n - 0| \ge \epsilon}
        \le \prob{X_n \neq 0}
        = \prob[\big]{U\le \tfrac1n}
        = \tfrac1n
        \to 0.
    \]

    \item[\ref{it: (conv in prob but not in Lp) not in L1}]
    We have
    \[
        \E{|X_n - 0|}
        = \E{X_n}
        = n \cdot \prob[\big]{U \in [0, 1/n)}
        = n \cdot \frac1n
        = 1 \not\to 0.
    \]            

    \item[\ref{it: (conv in prob but not in Lp) for all q}]
    Define \(Y_n \coloneq X_n^{1/p}\) with \(X_n\) as above. Then
    \[
        \E{|Y_n - 0|^q} = \E{X_n^{q/p}} = n^{q/p} \cdot \prob[\big]{U \in [0, 1/n)} = n^{q/p - 1}.
    \]
    This converges to \(0\) for all \(q < p\) but not for \(q \ge p\).
\end{enumerate}    
\end{solution}


\begin{exercise}[subtitle={Random vectors in high dimension}]
    In this exercise we will see that high-dimensional spaces have some counter-intuitive
    properties. Let \((X^n)_{n\in \nat}\) and \((Y^n)_{n\in \nat}\) be two independent
    sequences of independent identically distributed random variables with mean
    zero and finite variance \(\sigma_X^2\) and \(\sigma_Y^2\) respectively. Let \(X_\dims \coloneq (X^1, \ldots, X^\dims)\)
    and \(Y_\dims \coloneq (Y^1, \ldots, Y^\dims)\) be the corresponding random vectors in
    \(\real^\dims\).
    Prove that
    \begin{enumerate}[label=(\alph*), noitemsep]
        \item\label{it: (random vectors in high dim) norms}
        asymptotically and in relative terms both vectors have the same
        length in high dimension, that is
        \[
            \frac{\norm{X_\dims}^2}{\norm{Y_\dims}^2} \overset\as\to \frac{\sigma_X^2}{\sigma_Y^2} \qquad (\dims \to \infty).
        \]
        \item\label{it: (random vectors in high dim) angle} their cosine angle is zero asymptotically, that is
        \[
            \scp[\Big]{\frac{X_\dims}{\norm{X_\dims}}, \frac{Y_\dims}{\norm{Y_\dims}}} \overset\as\to 0 \qquad (\dims \to \infty).
        \]
        In other words: random vectors are asymptotically orthogonal in high dimension.
    \end{enumerate}
\end{exercise}

\begin{solution}
\begin{enumerate}[wide=0pt]
    \item[\ref{it: (random vectors in high dim) norms}]
    By definition we have
    \[
        \frac1\dims\norm{X_\dims}^2
        = \frac1\dims \sum_{i=1}^\dims (X^i)^2
        \overset{\as}\to \E{(X^i)^2}
        = \sigma_X^2
    \]
    By a similar argument we have \(\frac1\dims\norm{Y_\dims}^2\overset{\as}\to \sigma_Y^2\)
    and therefore by continuous mapping and joint convergence
    \[
        \frac{\norm{X_\dims}^2}{\norm{Y_\dims}^2}
        = \frac{\frac1\dims\norm{X_\dims}^2}{\frac1\dims\norm{Y_\dims}^2}
        \overset{\as}\to \frac{\sigma_X^2}{\sigma_Y^2}.
    \]
    \item[\ref{it: (random vectors in high dim) angle}]
    By continuous mapping we have \(\sqrt{\frac1\dims \norm{X_\dims}^2} \to \sigma_X\)
    and consequently
    \[
        \frac{\scp{X_\dims, Y_\dims}}{\norm{X_\dims} \norm{Y_\dims}}
        = \frac{\frac1\dims \sum_{i=1}^\dims X^i Y^i}{\sqrt{\frac1\dims \norm{X_\dims}^2} \sqrt{\frac1\dims \norm{Y_\dims}^2}}
        \overset{\as}\to \frac{\E{X^i Y^i}}{\sigma_X\sigma_Y}
        = 0.
    \]
\end{enumerate}
\end{solution}



\begin{exercise}[subtitle={Find the mistake}]
    Let \((X_n)_{n\in \nat}\) and \(X\) be random variables in \(\real^\dims\).
    Which step in the following chain of implications is \textbf{wrong}?
    \[
        X_n \overset{d}\to X
        \overset{?}\implies X_n - X \overset{d}\to 0
        \overset{?}\implies X_n - X \overset{p}\to 0
        \overset{?}\implies X_n \overset{p}\to X.
    \]
    Justify all other steps and find a counterexample for the wrong step.
\end{exercise}

\begin{solution}
    The first step is wrong. Consider \(X_n = -X\) with
    \(\prob{X=-1}=\prob{X=1}=\frac12\). Then
    as in  Exercise \ref{ex: convergence in distribution but not in
    probability} \(X_n \to X\), but \(X_n - X = -2X\) does not converge to
    zero in distribution.

    All the other steps are correct. The second step is the special case
    that convergence in distribution to a constant implies convergence in
    probability to that constant. The third step follows from 
    \[
        \prob{\norm{X_n - X} > \epsilon}
        =\prob[\big]{\norm{(X_n - X) - 0} > \epsilon}
        \to 0.
        \qedhere
    \]
\end{solution}



\begin{exercise}[subtitle={Uncorrelated but not independent normal}]
    \label{ex: uncorrelated but not independent normal}
    Let \(X \sim \normal{0,1}\) and \(\epsilon \sim \uniform(\set{-1,1})\), that is \(\prob{\epsilon = 1} = \prob{\epsilon = -1} = \frac{1}{2}\), be independent
    random variables and define
    \(Y \coloneq \epsilon X\).
    Show that \(X\) and \(Y\) are both normal distributed and uncorrelated but not independent.
\end{exercise}

\begin{solution}
    We have \(X \sim \normal{0,1}\) by definition and \(Y\) is normal
    distributed since it is a product of the independent random variables
    \(\epsilon\) and \(X\) and \(\epsilon\) is
    symmetric around zero. More formally
    \[\begin{aligned}
        \prob{Y \le y}
        &= \prob{\epsilon X \le y}
        \\
        &= \prob{X \le y \mid \epsilon = 1} \prob{\epsilon = 1} + \prob{-X \le y \mid \epsilon = -1} \prob{\epsilon = -1}
        \\
        &= \frac{1}{2} \prob{X \le y} + \frac{1}{2} \prob{-X \le y} 
        \\
        &= \prob{X \le y}.
    \end{aligned}\]
    where we used the symmetry of \(X\) in the last step. Since this is true for all bounded continuous functions \(g\) we have \(Y \overset{d}= X\).
    We have
    \[
        \cov{X}{Y}
        = \E{XY} - \E{X}\E{Y}
        = \E{\epsilon X^2} - \E{\epsilon}\E{X}^2
        = 0 - 0 = 0,
    \]
    so \(X\) and \(Y\) are uncorrelated.
    However, they are not independent since
    \[
        \prob{\abs{X} > 1, \abs{Y} < 1}
        = \prob{\abs{Y} >1, \abs{Y} < 1} = 0
        < \prob{\abs{X} > 1}\prob{\abs{Y} < 1}.
        \qedhere
    \]
\end{solution}



\begin{exercise}[subtitle={Infinite Monkey Theorem}]
    The infinite monkey theorem states that an immortal monkey hitting keys at
    random on a typewriter keyboard for an infinite amount of time will
    eventually type any arbitrarily long text -- for example all the
    works of William Shakespeare -- with probability one.  It is also known as
    the ``law of truly large numbers''.
    Prove it!
    \begin{enumerate}[label={\emph{Level \arabic*:}},ref={Level \arabic*},leftmargin=4em,noitemsep]
        \item Assume the hit keys are independent and uniformly distributed.
        \item\label{it: level iid} Assume the hit keys are iid with a probability distribution that
        assigns each of the (finite) keys a positive probability.
    \end{enumerate}
\end{exercise}

\begin{solution}
    Let \(\alphabet\) be a finite alphabet of keys and let \(X_n\in
    \alphabet\) be the letter typed at time \(n\). Let \(y\in \alphabet^m\)
    be the desired text of length \(m\). We essentially assume \(m=1\)
    without loss of generality with the definition of
    \[
        Y_n = (X_{nm}, X_{nm+1}, \ldots, X_{m(n+1)-1}) \in \alphabet^m,
    \]
    which is still a sequence of iid random variables in \ref{it: level iid}
    (or still a Markov chain in \ref{it: level markov chain}).  Let us consider \ref{it: level iid}
    first. Since
    \[\begin{aligned}
        \delta \coloneq \prob{Y_n = y}
        &= \prob{X_{nm} = y_1, X_{nm+1} = y_2, \ldots, X_{m(n+1)-1} = y_m}
        \\
        &= \prod_{i=1}^m \prob{X_{nm+i-1} = y_i} > 0,
    \end{aligned}
    \]
    Due to \(\sum_{n=1}^\infty \prob{Y_n = y} = \sum_{n=1}^\infty \delta = \infty\)
    we may apply the second Borel-Cantelli lemma \eqref{eq: borel cantelli 2} to get
    \[
        \prob[\big]{Y_n = y \text{ for infinitely many } n}
        = \prob[\big]{\limsup_{n\to\infty} \set{Y_n = y}}
        \overset{\text{\eqref{eq: borel cantelli 2}}}= 1.
    \]
    Observe that we have proven more than the desired claim: The monkey will
    not only type the desired text \(y\) at least once but infinitely many
    times with probability one!
\end{solution}



\begin{exercise}[subtitle={CGF and convex conjugate calculations}]
    \label{ex: convex conjugate calculations}
    Prove the formulas for CGF and their convex conjugates in the following
    table
    \begin{center}
    \def\arraystretch{1.2}
    \begin{tabular}{l l l}
        \toprule
        \(\pSym_X\)
        &  \(\CGF_X(t)\) & \(\CGF_X^*(a)\)
        \\
        \midrule
        \(\Poisson(\lambda)\)
        & \(\lambda (e^{t} - 1)\)
        & \(1-a + \ln(\frac{a}\lambda)a\) 
        \\
        \(\exponential(\lambda)\)
        & \(\ln(\frac{\lambda}{\lambda - t})\) for \(t < \lambda\)
        & \(\lambda a -1 - \ln(\lambda a)\)
        \\
        \(\bernoulli(p)\)
        & \(\ln(1 - p + p e^{t})\)
        & \(a\ln(\frac{a}p) + (1-a)\ln(\frac{1-a}{1-p})\)
        \\
        \bottomrule
    \end{tabular}
    \end{center}
\end{exercise}
\begin{solution}
    The formulas for the CGF follow directly from the formulas for the MGF (Exercise \ref{ex: compute mgf}).
    Since the cumulant generating function is always convex (Lemma \ref{lem: properties of CGF} \ref{it: CGF convex}),
    the supremum in the definition of the convex conjugate
    \[
        \CGF_X^*(a) = \sup_{t \in \real} \{t a - \CGF_X(t)\}
    \]
    can be found using the first order condition
    \[
        a = \CGF_X'(t).
    \]
    Recall: the convex conjugate finds the point \(t\) for the gradient \(a\).
    The derivatives and maximizers \(t^*\) from the first order condition are thus given by
    \begin{center}
    \def\arraystretch{1.2}
    \begin{tabular}{l l l l}
        \toprule
        \(\pSym_X\)
        & \(\CGF_X(t)\)
        & \(\CGF_X'(t)\)
        & \(t^*=\argmax_t (ta - \CGF_X(t))\)
        \\
        \midrule
        \(\Poisson(\lambda)\)
        & \(\lambda (e^{t} - 1)\)
        & \(\lambda e^t\)
        & \(\ln\bigl(\frac{a}{\lambda}\bigr)\)
        \\
        \(\exponential(\lambda)\)
        & \(\ln(\frac{\lambda}{\lambda - t})\) for \(t < \lambda\)
        & \(\frac1{\lambda-t}\) for \(t<\lambda\)
        & \(\lambda - \frac1a\)
        \\
        \(\bernoulli(p)\)
        & \(\ln(1 - p + p e^{t})\)
        & \(\frac{pe^t}{1-p+p e^t}\)
        & \(\ln(\frac{a}p\frac{1-p}{1-a})\)
        \\
        \bottomrule
    \end{tabular}
    \end{center}
    where \(t^*\) of the Bernoulli distribution follows from
    \begin{align*}
        a \overset{!}= \frac{pe^t}{1-p+p e^t}
        &\iff a(1-p+p e^t) = pe^t
        \\
        &\iff a(1-p) = pe^t(1 - a)
        \iff t = \ln(\tfrac{a}p\tfrac{1-p}{1-a}).
    \end{align*}
    Plugging the maximizers \(t^*\) into \(ta - \CGF_X(t)\) yields the desired formulas.
    For example for the Bernoulli distribution, we have
    \[
        \CGF_X^*(a) = a \ln(\tfrac{a}p\tfrac{1-p}{1-a}) - \ln(\underbrace{1 - p + p \tfrac{a}p\tfrac{1-p}{1-a}}_{= (1-p)(1+\frac{a}{1-a})\mathrlap{=\frac{1-p}{1-a}}})
        = a\ln(\tfrac{a}p) + (1-a)\ln(\tfrac{1-a}{1-p}).
        \qedhere
    \]
\end{solution}


% 
\begin{exercise}[subtitle={Randomized linear algebra \Coffeecup}]
    Let \(x,y \in \real^\dims\), where the dimension \(\dims\) is really large.
    The computation of \(\scp{x,y}\) clearly costs \(\dims\) multiplications and
    additions. We will try to speed it up using a random sample.
    Specifically, let \(Z_1, \ldots Z_n\) be independently drawn from
    \(\frac1\dims\sum_{i=1}^\dims \delta_{x_iy_i}\), that is
    \(\prob{Z_1 = x_iy_i} = \frac1\dims\) for \(i=1, \ldots, \dims\).
    \begin{enumerate}[label=(\alph*), noitemsep]
        \item\label{it: (randomized linear algebra) expectation}
        Prove that \(\dims \E{Z_1} = \scp{x,y}\),

        \item\label{it: (randomized linear algebra) convergence}
        Assuming \(\scp{x,y} \neq 0\), prove the following estimator
        \[
            \hat S_n \coloneq \frac\dims{n} \sum_{i=1}^n Z_i,
        \]
        converges against the scalar product in relative terms, that is
        \[
            \frac{\hat S_n}{\scp{x,y}} \overset\as\to 1 \qquad (n \to \infty).
        \]

        \item\label{it: (randomized linear algebra) PAC bound}
        For \(\hat S_n\) to be an effective speedup of the
        computation of the scalar product it is necessary that \(n\le \dims\). So
        clearly \(n\to \infty\) is not an interesting statement.
        Let \(m = \max_i \abs{x_i y_i}\) and show that
        \[
            \prob[\Big]{\abs[\Big]{\frac{\hat S_n}{\scp{x,y}} - 1} > \epsilon}
            \le 2\exp\Bigl(-n \epsilon^2 \frac{\E{Z_1}^2}{8 m^2}\Bigr).
        \]
        \textbf{Hint:} Corollary \ref{cor: hoeffding inequality for bounded random variables}.
        \begin{remark}[Many significant values]
            If most \(x_iy_i\) are very small except for a few \(i=\set{1,\dots, \dims}\),
            then the inner product \(\scp{x,y}\) is dominated by these values.
            This means that a random sample can be very far off if it does not
            contain these significant indices. This case is captured
            by the ratio \(\frac{\E{Z_1}}{m}\): If the average product
            \(\E{Z_1}\) is much smaller, than largest product \(m\), the sum is
            dominated by a few large entries and convergence is slow.
        \end{remark}

        \item\label{it: (randomized linear algebra) sample complexity}
        Assume \(\frac{\E{Z_1}}{m} \ge 0.1\), that a precision of \(\epsilon
        = 10^{-4}\) is required and a \(99\%\) confidence. How big must
        \(\dims\) be such that this approach reduces the computational cost?
        How big to get a \(99.999\%\) confidence?
    \end{enumerate}
\end{exercise}

\begin{solution}
\begin{enumerate}[wide=0pt]
    \item[\ref{it: (randomized linear algebra) expectation}]
    By definition we have
    \[
        \E{Z_1} = \frac1\dims\sum_{i=1}^\dims  x_i y_i = \frac1\dims \scp{x,y}.
    \]            

    \item[\ref{it: (randomized linear algebra) convergence}]
    We have \(\frac{\hat S_n}{\scp{x,y}} = \frac1n \sum_{i=1}^n Z_i /
    \E{Z_1}\) and thus the claim by the strong law of large numbers.

    \item[\ref{it: (randomized linear algebra) PAC bound}]
    Since \(Z_i \in [-m,m]\) for \(i=1, \ldots, n\) we can apply
    Hoeffding's inequality for bounded random variables (Corollary \ref{cor:
    hoeffding inequality for bounded random variables}) to get
    \begin{align*}
        \prob[\Big]{\abs[\Big]{\frac{\hat S_n}{\scp{x,y}} - 1} > \epsilon}
        &=\prob[\Big]{\abs[\Big]{\frac1n\sum_{i=1}^n Z_i - \E{Z_1}} > \epsilon \abs{\E{Z_1}}}
        \\
        &\le 2\exp\Bigl(-n \epsilon^2 \frac{\E{Z_1}^2}{2(2m)^2}\Bigr).
    \end{align*}

    \item[\ref{it: (randomized linear algebra) sample complexity}]
    We need \(n\) to satisfy
    \[
        2\exp\Bigl(-n \epsilon^2 \frac{\E{X_1 Y_1}^2}{8 m^2}\Bigr) \le \delta 
        \iff n \ge \frac{8\ln(\frac1{2\delta})}{\epsilon^2 \frac{\E{X_1 Y_1}^2}{m^2}}.
    \]
    Since \(\E{X_1 Y_1}^2 \ge 0.1\) we have that
    \[
        n \ge \frac{80\ln(\frac1{2\delta})}{\epsilon^2}
    \]
    is sufficient to get the required confidence. We now simply
    plug in the numbers. For \(\delta = 0.01 = 1\%\) we have
    \[
        n \ge \frac{80\ln(\frac1{0.02})}{10^{-8}} \approx 3.12 \cdot 10^{10}.
    \]
    For \(\delta = 0.00001 = 0.001\%\) we have
    \[
        n \ge \frac{80\ln(\frac1{0.00002})}{10^{-8}} \approx 8.66 \cdot 10^{10}.
    \]
    If \(n\le \dims\) then this approach reduces the computational cost.
    So for \(\dims \ge 3.12 \cdot 10^{10}\) (or \(\dims \ge 8.66 \cdot 10^{10}\) respectively)
    this approach reduces the computational cost. As in Remark \ref{rem:
    accuracy vs confidence} the confidence has almost no effect on the
    required dimension. In contrast the precision \(\epsilon\)
    dominates the required dimension.
    \qedhere
\end{enumerate}
\end{solution}

% 
\begin{exercise}[subtitle={Properties of the cumulant generating function}]
    \label{ex: properties cumulant generating function}
    Let \(X\) be a random variable. Then the cumulant generating function
    \(\CGF_X(t) = \log \MGF_X(t)\) has the following properties:
    \begin{enumerate}[label=(\roman*), noitemsep]
        \item\label{it: cgf 0=0} \(\CGF_X(0) = 0\).
        \item\label{it: cgf linear} \(\CGF_{aX + b}(t) = \scp{b, t} + \CGF_X(\scp{a, t})\) for any \(a \in \real\), \(b,t\in \hilbert\).
        \item\label{it: cumulant independent sums} Let \(Y\) be \underline{independent} from \(X\), then
        \[
            \CGF_{X+Y}(t) = \CGF_X(t) + \CGF_Y(t)
        \]
    \end{enumerate}
    If the moment generating function is finite in a neighborhood of zero,
    then \(\CGF_X\) is infinitely often differentiable and
    \begin{enumerate}[label=(\roman*), resume, noitemsep]
        \item \(\CGF_X'(0) = \E{X}\),
        \item \(\CGF_X''(0) = \var{X}\)
    \end{enumerate}
\end{exercise}

\begin{solution}
    For \ref{it: cgf 0=0}-\ref{it: cumulant independent sums} follow from the
    corresponding properties of the MGF (Exercise \ref{ex: properties moment
    generating functions}) by application of the logarithm.  For example,
    \(\CGF_X(0)=0\) follows directly from \(\MGF_X(0)=1\). For convexity
    let \(t, s\in \hilbert\) and \(\lambda \in [0,1]\). 

    The derivatives are straightforward calculations using the derivatives of the MGF
    and the chain rule due to \(\CGF_X(t) = \log \MGF_X(t)\). That is
    \begin{alignat*}{3}
        \CGF_X'(t)
        &= \frac{\MGF_X'(t)}{\MGF_X(t)}
        \qquad &\implies &\CGF_X'(0) = \E{X},
        \\
        \CGF_X''(t)
        &= \frac{\MGF_X''(t)}{\MGF_X(t)} - \frac{(\MGF_X'(t))^2}{(\MGF_X(t))^2}
        \qquad &\implies &\CGF_X''(0)
        \begin{aligned}[t]
        &= \E{X^2} - (\E{X})^2
        \\
        &= \var{X}.
        \end{aligned}
    \end{alignat*}
    This proves the claims.
\end{solution}


\end{document} 

%%% Local Variables: 
%%% mode: latex
%%% TeX-master: t
%%% End: 

